\begin{proof}[Proof of \cref{obj:lem:noise-cost-scaling}]
For \(0\le \sigma\le1\), write
\[
F_\sigma(h)=1+\Psi(h,\sigma),\qquad 0<h\le1.
\]
\begin{enumerate}
\item \emph{Continuity.}
Fix \(\sigma\in[0,1]\). If \(\sigma=0\), then
\[
\Psi(h,0)=0,\qquad 0<h\le1,
\]
so continuity follows. Suppose \(\sigma>0\). The power and logarithmic expressions defining the cost are continuous for positive \(h\): their denominators are positive, and the logarithm has positive argument. At their interface,
\[
\sqrt{\sigma^4}=\sigma^2,\qquad
\left(\frac{\sigma}{\sigma^4}\right)^{2/3}
=\sigma^{-2}
=(\sigma^4)^{-1/2}
 \left[1+\log\left(\frac{\sigma^2}{\sqrt{\sigma^4}}\right)\right].
\]
Thus the two expressions for the cost agree at \(h=\sigma^4\).
% lean: CausalSmith.Stat.RdTruesideNoiseFrontier.noiseCost_compact_interface
At the direct-branch interface, \(\sigma^4\le\sigma\), since \(0<\sigma\le1\), and
\[
\left(\frac{\sigma}{\sigma}\right)^{2/3}-1=0.
\]
Gluing the continuous expressions at these interfaces proves continuity on \((0,1]\), including when \(\sigma=1\) and the interfaces coincide.
% lean: CausalSmith.Stat.RdTruesideNoiseFrontier.noiseCost_continuousOn

\item \emph{Nonnegativity.}
Fix \(0<h\le1\) and \(0\le\sigma\le1\). If \(\sigma\le h\), the cost is zero. If \(h<\sigma\) and \(\sigma^4\le h\), then
\[
\frac{\sigma}{h}\ge1,\qquad
\left(\frac{\sigma}{h}\right)^{2/3}\ge1,
\]
so the power-branch cost is nonnegative. In the remaining case, \(h<\sigma^4\), positivity of \(h\) gives
\[
\sqrt h\le\sigma^2,\qquad
\log\left(\frac{\sigma^2}{\sqrt h}\right)\ge0,
\qquad h^{-1/2}\ge1.
\]
Consequently,
\[
h^{-1/2}\left[1+\log\left(\frac{\sigma^2}{\sqrt h}\right)\right]\ge1,
\]
which proves nonnegativity on the logarithmic branch as well.
% lean: CausalSmith.Stat.RdTruesideNoiseFrontier.noiseCost_nonneg

\item \emph{Monotonicity.}
Fix \(0<h\le t\le1\). If \(\sigma\le h\), both costs are zero. If \(h<\sigma\le t\), nonnegativity gives
\[
\Psi(t,\sigma)=0\le\Psi(h,\sigma).
\]
It remains to consider \(h\le t<\sigma\), so \(\sigma>0\).

If \(\sigma^4\le h\), both resolutions lie on the power branch, and
\[
0\le\frac{\sigma}{t}\le\frac{\sigma}{h}
\quad\Longrightarrow\quad
\left(\frac{\sigma}{t}\right)^{2/3}-1
\le
\left(\frac{\sigma}{h}\right)^{2/3}-1.
\]
% lean: CausalSmith.Stat.RdTruesideNoiseFrontier.noiseCost_power_antitone

If \(h<\sigma^4\le t\), compare through the interface:
\[
\begin{aligned}
\Psi(t,\sigma)
&=\left(\frac{\sigma}{t}\right)^{2/3}-1\\
&\le\left(\frac{\sigma}{\sigma^4}\right)^{2/3}-1\\
&=(\sigma^4)^{-1/2}
 \left[1+\log\left(\frac{\sigma^2}{\sqrt{\sigma^4}}\right)\right]-1\\
&\le h^{-1/2}
 \left[1+\log\left(\frac{\sigma^2}{\sqrt h}\right)\right]-1
=\Psi(h,\sigma).
\end{aligned}
\]
The first inequality is the power comparison just established, and the equality is the interface identity. For the last inequality,
\[
(\sigma^4)^{-1/2}\le h^{-1/2},
\qquad
0=\log\left(\frac{\sigma^2}{\sqrt{\sigma^4}}\right)
\le\log\left(\frac{\sigma^2}{\sqrt h}\right);
\]
multiplying the corresponding nonnegative factors gives the required bound.
% lean: CausalSmith.Stat.RdTruesideNoiseFrontier.noiseCost_compact_antitone

Finally, if \(h\le t<\sigma^4\), then
\[
\sqrt t\le\sigma^2,\qquad
0\le\log\left(\frac{\sigma^2}{\sqrt t}\right)
\le\log\left(\frac{\sigma^2}{\sqrt h}\right),
\qquad
t^{-1/2}\le h^{-1/2}.
\]
Here the logarithmic comparison follows from \(\sqrt h\le\sqrt t\) and monotonicity of the logarithm on positive arguments. Multiplication yields
\[
t^{-1/2}\left[1+\log\left(\frac{\sigma^2}{\sqrt t}\right)\right]-1
\le
h^{-1/2}\left[1+\log\left(\frac{\sigma^2}{\sqrt h}\right)\right]-1.
\]
These cases establish that the cost is nonincreasing.
% lean: CausalSmith.Stat.RdTruesideNoiseFrontier.noiseCost_antitoneOn

\item \emph{Ratio bounds.}
Assume \(0<h\le t\le\sigma\le1\). Then \(t,\sigma>0\), and nonnegativity gives
\[
F_\sigma(t)\ge1>0.
\]
We first prove the multiplicative comparisons
\[
\left(\frac th\right)^{1/2}F_\sigma(t)
\le F_\sigma(h)
\le\frac th F_\sigma(t).
\]

If \(\sigma^4\le h\), the power formula holds at both resolutions:
\[
F_\sigma(h)=\left(\frac{\sigma}{h}\right)^{2/3},
\qquad
F_\sigma(t)=\left(\frac{\sigma}{t}\right)^{2/3}.
\]
These identities include \(t=\sigma\), since the direct-branch value there is one. Hence
\[
F_\sigma(h)=\left(\frac th\right)^{2/3}F_\sigma(t).
\]
Because \(t/h\ge1\),
\[
\left(\frac th\right)^{1/2}
\le\left(\frac th\right)^{2/3}
\le\frac th.
\]
Multiplying by \(F_\sigma(t)>0\) proves both comparisons in this case.
% lean: CausalSmith.Stat.RdTruesideNoiseFrontier.noiseCost_power_comparison

Suppose next that \(h<\sigma^4\) and \(t\le\sigma^4\). The logarithmic formula, extended to the interface by the identity in the first step, gives
\[
F_\sigma(h)=h^{-1/2}
 \left[1+\log\left(\frac{\sigma^2}{\sqrt h}\right)\right],
\qquad
F_\sigma(t)=t^{-1/2}
 \left[1+\log\left(\frac{\sigma^2}{\sqrt t}\right)\right].
\]
The same formulas hold whenever \(0<h\le t\le\sigma^4\).
For such resolutions define
\[
\eta_{\mathrm{cmp}}=\frac{\sqrt t}{\sqrt h},
\qquad
\Lambda_{\mathrm{cmp}}=\log\left(\frac{\sigma^2}{\sqrt t}\right).
\]
Then
\[
\eta_{\mathrm{cmp}}\ge1,\qquad
\Lambda_{\mathrm{cmp}}\ge0,\qquad
\log\left(\frac{\sigma^2}{\sqrt h}\right)
=\Lambda_{\mathrm{cmp}}+\log\eta_{\mathrm{cmp}}.
\]
Using \(0\le\log\eta_{\mathrm{cmp}}\le\eta_{\mathrm{cmp}}-1\), we obtain
\[
\begin{aligned}
1+\Lambda_{\mathrm{cmp}}
&\le1+\Lambda_{\mathrm{cmp}}+\log\eta_{\mathrm{cmp}},\\
1+\Lambda_{\mathrm{cmp}}+\log\eta_{\mathrm{cmp}}
&\le\Lambda_{\mathrm{cmp}}+\eta_{\mathrm{cmp}}
\le\eta_{\mathrm{cmp}}(1+\Lambda_{\mathrm{cmp}}).
\end{aligned}
\]
The final inequality uses
\((\eta_{\mathrm{cmp}}-1)\Lambda_{\mathrm{cmp}}\ge0\).
Since
\[
h^{-1/2}=\frac1{\sqrt h},\qquad
t^{-1/2}=\frac1{\sqrt t},\qquad
\eta_{\mathrm{cmp}}=\left(\frac th\right)^{1/2},
\qquad
\eta_{\mathrm{cmp}}^2=\frac th,
\]
multiplication by \(1/\sqrt h\) gives
\[
\left(\frac th\right)^{1/2}F_\sigma(t)
\le F_\sigma(h)
\le\frac th F_\sigma(t).
\]
% lean: CausalSmith.Stat.RdTruesideNoiseFrontier.noiseCost_compact_comparison

The remaining case is \(h<\sigma^4<t\). Apply the logarithmic comparison from \(h\) to \(\sigma^4\), and the power comparison from \(\sigma^4\) to \(t\). As \(\sigma^4>0\), their lower bounds compose as
\[
\begin{aligned}
\left(\frac th\right)^{1/2}F_\sigma(t)
&=\left(\frac{\sigma^4}{h}\right)^{1/2}
  \left(\frac{t}{\sigma^4}\right)^{1/2}F_\sigma(t)\\
&\le\left(\frac{\sigma^4}{h}\right)^{1/2}F_\sigma(\sigma^4)
\le F_\sigma(h),
\end{aligned}
\]
and their upper bounds compose as
\[
F_\sigma(h)
\le\frac{\sigma^4}{h}F_\sigma(\sigma^4)
\le\frac{\sigma^4}{h}\frac{t}{\sigma^4}F_\sigma(t)
=\frac th F_\sigma(t).
\]
% lean: CausalSmith.Stat.RdTruesideNoiseFrontier.noiseCost_comparison_trans
Dividing the comparisons by the positive quantity \(F_\sigma(t)\) proves
\[
\left(\frac th\right)^{1/2}
\le\frac{F_\sigma(h)}{F_\sigma(t)}
\le\frac th.
\]
% lean: CausalSmith.Stat.RdTruesideNoiseFrontier.noiseCost_ratio_bounds

\item \emph{Dilation bound.}
Let \(0<h\le1\), \(0\le\sigma\le1\), \(a\ge1\), and \(ah\le1\). If \(\sigma\le ah\), then
\[
F_\sigma(ah)=1
\le\max\{1,a^{-1/2}F_\sigma(h)\}.
\]
This case includes \(\sigma=0\). Otherwise, \(ah<\sigma\). Since \(a>0\),
\[
0<h\le ah\le1,\qquad F_\sigma(ah)>0.
\]
Applying the lower ratio bound with \(t=ah\) gives
\[
a^{1/2}
\le\frac{F_\sigma(h)}{F_\sigma(ah)},
\qquad
a^{1/2}F_\sigma(ah)\le F_\sigma(h).
\]
Multiplication by \(a^{-1/2}\ge0\), together with
\(a^{-1/2}a^{1/2}=1\), yields
\[
F_\sigma(ah)
\le a^{-1/2}F_\sigma(h)
\le\max\{1,a^{-1/2}F_\sigma(h)\}.
\]
Substituting \(F_\sigma(h)=1+\Psi(h,\sigma)\) proves the last assertion.
% lean: CausalSmith.Stat.RdTruesideNoiseFrontier.noiseCost_dilation_of_lower_ratio
\end{enumerate}
\end{proof}
